% Аннотации докладов научного семинара
% Лаборатория математики ФН1
% Компиляция: pdflatex annotacii-dokladov.tex

\documentclass[12pt, a4paper]{article}

\usepackage[T2A]{fontenc}
\usepackage[utf8]{inputenc}
\usepackage[russian]{babel}
\usepackage[left=22mm, right=20mm, top=20mm, bottom=20mm]{geometry}
\usepackage{amsmath, amssymb}

\pagestyle{plain}
\setlength{\parindent}{0pt}
\setlength{\parskip}{0.4em}

\newcommand{\seminar}[4]{%
  \vspace{0.9em}
  {\bfseries #1} \quad {\small\itshape #2}\\
  {\large #3}\\
  {\small Докладчик: #4}\par
  \vspace{0.2em}
}

\begin{document}

\begin{center}
  {\Large\bfseries Аннотации докладов}

  \vspace{0.3em}
  {\small Научный семинар Лаборатории математики ФН1 · осенний семестр 2026/27\\
  Аудитория 312, среда, 16:30}
\end{center}

\vspace{1em}

\hrule

\seminar{9 сентября}{Другое}{Открытие семестра: план работы семинара}{руководитель семинара}
Обсуждаем тематику осеннего семестра, распределяем даты докладов и назначаем
оппонентов. Приглашаются все желающие выступить: заявку можно подать прямо
на заседании.

\hrule

\seminar{16 сентября}{Учебная лекция}{Асимптотические методы в задачах теории возмущений}{проф. кафедры}
Обзорная лекция для студентов старших курсов и аспирантов. Разбираются метод
малого параметра, метод сращиваемых асимптотических разложений и типичные
ошибки при построении асимптотики. Предполагается знание курса
дифференциальных уравнений.

\hrule

\seminar{23 сентября}{Разбор статьи}{Устойчивость разностных схем для уравнения переноса}{аспирант 2-го года}
Ридинг-клуб: совместное чтение современной публикации. Текст статьи
рассылается участникам за неделю до заседания. Разбираем доказательство
основной теоремы и обсуждаем, насколько условие устойчивости необходимо.

\hrule

\seminar{30 сентября}{Научное сообщение}{Спектральные свойства оператора Штурма~--- Лиувилля с разрывным коэффициентом}{доц. кафедры}
Излагаются результаты собственного исследования: получены оценки роста
собственных значений и доказана полнота системы собственных функций
при ослабленных условиях на коэффициент.

\hrule

\seminar{7 октября}{Практикум}{Символьные вычисления в SymPy для проверки выкладок}{доц. кафедры}
Занятие за компьютерами. Показываем, как проверять производные, интегралы,
разложения в ряд и решения дифференциальных уравнений. Обсуждаем, где
символьные пакеты ошибаются и почему нельзя доверять им безоговорочно.
Ноутбук с примерами выкладывается после занятия.

\hrule

\seminar{14 октября}{Учебная лекция}{Методы оптимизации в задачах машинного обучения: взгляд математика}{приглашённый докладчик}
Что из классической теории оптимизации действительно работает в обучении
нейронных сетей, а что нет. Сходимость стохастического градиентного спуска,
роль условия Липшица, современные адаптивные методы.

\hrule

\seminar{21 октября}{Разбор статьи}{Априорные оценки для эллиптических операторов}{магистранты}
Второй ридинг-клуб семестра. Доклад готовят двое магистрантов:
один излагает постановку и результат, второй~--- технику доказательства.

\hrule

\seminar{11 ноября}{Научное сообщение}{Численное решение обратных задач теплопроводности}{аспирант 1-го года}
Первое выступление докладчика на семинаре. Рассматривается задача
восстановления начального распределения температуры по измерениям
в конечный момент времени, обсуждается регуляризация по Тихонову.

\hrule

\seminar{25 ноября}{Практикум}{Подготовка публикации в LaTeX и работа с библиографией}{доц. кафедры}
Разбираем структуру статьи, оформление формул и таблиц, работу с BibTeX,
подготовку материалов к отправке в журнал. Полезно всем, кто готовит
курсовую, диплом или первую статью.

\hrule

\seminar{9 декабря}{Другое}{Короткие сообщения участников: итоги семестра}{открытая сессия}
Заключительное заседание. Каждый участник рассказывает о своей работе
за 10 минут. Заявки принимаются до 1 декабря.

\hrule

\vspace{1.2em}

{\small
\textbf{Как заявить доклад.} Обсудите тему с научным руководителем и направьте
секретарю семинара аннотацию на 5--7 предложений с желаемой датой. Регламент:
основной доклад~--- 40 минут, короткое сообщение~--- 10 минут. Презентация
принимается в форматах PDF или Beamer.
}

\end{document}
